\documentclass{article}
\usepackage[T1]{fontenc}
\usepackage{lmodern}
\usepackage{graphicx}
\usepackage{amsmath,amssymb}
\usepackage[a4paper,margin=2cm]{geometry}
\usepackage[absolute,overlay]{textpos}
\setlength{\TPHorizModule}{1mm}
\setlength{\TPVertModule}{1mm}
\usepackage[indonesian]{babel}
\usepackage{hyperref}
\setlength{\emergencystretch}{2em}
% unit: Halaman 5

\begin{document}

% === Halaman 5 (python/native) ===

5

Bit, Byte, dan Kode Heksadesimal 

• Pesan di dalam cipher modern dienkripsi bit-per-bit atau byte-per-byte, atau dalam kelompok bit (byte). 


1 byte = 8 bit

• Pada beberapa algoritma kriptografi, pesan dikodekan ke dalam kode heksadesimal (Hex). 


1 kode hex = 4 bit


0000 = 0  

0001 = 1  

0010 = 2  

0011 = 3


0100 = 4  

0101 = 5  

0110 = 6  

0111 = 7


1000 = 8  

1001 = 9  

1010 = A 

1011 = B


1100 = C  

1101 = D 

1110 = E  

1111 = F

• Contoh: Pesan 1001110101101010 dalam kode Hex dengan  cara membagi pesan menjadi blok 4-bit: 


1001  1101  0110  1010 = 9D6A

\end{document}
